% ADD SKLEARN STUFF (SEE JEDISCIKIT.PY)
% INCLUDING ONE-HOT ENCODING


\chapter{ML classifiers: Na\"{i}ve Bayes}

For some folks, the coolest and most exciting parts of Data Science are the
\textbf{machine learning (ML)} algorithms. These are the under-the-hood
techniques for doing prediction and analytics. We'll look at a few of them
briefly in this class, and you'll delve into them more fully in later courses.

Recall that a \textbf{classifier} is an algorithm/program to predict which of
several \textbf{labels} (or \textbf{classes}) to assign to a data point.
Perhaps we're trying to predict what political party someone is affiliated
with, and our labels are ``Democrat,'' ``Republican,'' and ``Green.'' Perhaps
we're trying to predict whether someone will default on a loan we're
considering giving them, and the labels are ``safe'' and ``high-risk.''

There are many different algorithms for classification, but they all use a set
of \textbf{labeled examples} as ``\textbf{training data}.'' The idea is that by
scrutinizing a bunch of examples for which we \textit{know} the label, we can
intelligently figure out what label to assign to new, \textbf{unlabeled
examples}.

The attributes of our training data (and new data) that we'll use for
prediction are sometimes called \textbf{features}. If our training data is in
a DataFrame, the features are essentially all the relevant columns except for
one, which is the \textbf{target}. Our goal, given a new data point that has
values for all the features, is to label its target with the label it's most
likely to have, assuming the training data are representative examples. 

A classifier's \textbf{performance} can be measured in several ways, but the
most straightforward is to score how many unlabeled examples it gives the
correct label to. This of course begs the question: if the example is
\textit{un}labeled, how can we know whether or not the classifier's label was
correct? What we have to do in practice is set aside some of our (precious)
labeled data as a set of \textbf{test data}. We sort of pretend we don't know
the answers for the test data, feed them to the algorithm which has been
trained on the \textit{training} points, and then score how well it did.

One of the simplest (and quickest) classifiers around is called the
\textbf{Na\"{i}ve Bayes classifier}, for reasons which will become apparent. In
some ways, it's the rock-bottom, no-frills default algorithm for
classification. One reason it's important is that it's simple and actually
works surprisingly (even embarrassingly) well in many cases. Another reason is
that even when it doesn't turn out to have the best performance,
it's still a \textbf{baseline} which other algorithms' performance can be
compared to. A researcher might report: ``our whiz-bang algorithm X got 89.4\%
of examples correct on this test set, compared with 83.7\% for straightforward
Na\"{i}ve Bayes.''

\section{Background: Probability}

Probability theory is an area of mathematics that deals with
\textit{quantifying uncertainty}. We may not be sure that something is going
to happen, but we want to be precise about \textit{how} sure we are, given
everything else we may know.

\subsection{Terms}

Some terms: an \textbf{outcome} is something that may (or may not) happen.
This varies widely depending on the domain: a child that's born may or may not
be male; a train passenger may or may not have red hair; a customer may or may
not buy a particular product. ``Male,'' ``red,'' and ``The Girl on the Train,
hardback edition'' are all outcomes in these examples.

An \textbf{event}, by contrast, is a \textit{set} (group) of outcomes that
we're interested in. Often we don't care so much whether or not the customer
specifically bought ``The Girl on the Train, hardback edition'' as whether they
bought \textit{any} book by Paula Hawkins, or maybe any hardback book, or
maybe any book at all (as opposed to a DVD).

Sometimes we'll use the symbol $\Omega$ (``omega'') to refer to the
\textbf{sample space}, which is simply the set of \textit{all possible
outcomes}. This will vary depending on what's under consideration; it might be
\{ male, female \} in one case and \{ red, blonde, brunette, bald, other \} in
another case.

We'll also use the notation $\overline{\phantom{xx}}$ (overbar) to mean ``any
outcome \textit{except} something.'' For instance, $\overline{\textrm{red}}$
means ``the train passenger had a hair color \textit{other than} red.''

\subsection{Probability measures}

A \textbf{probability measure} is a function that tells us intuitively how
likely an outcome is to occur. We'll use the simple notation ``P($\cdot$)'' for
this, read as ``the probability of (something).'' The probability of any
outcome is always \textit{a number between 0 and 1}. 0 means it can never
possibly occur, 1 means it will definitely occur, and there's a continuum in
between. In our first example, we might say P(male) = .5, P(red) = .1, and
P(The Girl on the Train, hardback edition) = .0001.

You can talk about the probability of an event too, of course, and it's simply
the sum of the outcome probabilities. P(male-or-female) = 1, P(red or blonde
or bald) = .6, P(any hardcover book) = .08. We'll say that P($\Omega$) will be
1: after all, \textit{something} has to happen, even if it's ``the customer
didn't buy anything,'' an outcome to which we'll assign a probability.

\subsection{``Joint'' probability}

Very commonly, we'll be considering the likelihood that two different events
\textit{both} occur. For instance, perhaps we're interested in the probability
that a train passenger is a red-headed female. We'll write this as P(red,
female). The comma signals \textbf{joint} probability, which simply means that
both of these outcomes/events must occur in order for the combined event to be
considered to have occurred. If we wrote P(red, $\overline{\textrm{female}}$),
it would mean `the probability that the passenger is a red-headed
\textit{non}-female.

If you think about it, you'll realize that the joint probability cannot be
greater than either of the two individual probabilities. Adding another
condition (not female) to the first (red-headed) can't possibly
\textit{increase} the chance of it happening. Other than that, we don't know
how to compute the joint probability from the individual ones, though (yet).

Finally, note that P(A,B) is the same thing as P(B,A).

\subsection{``Conditional'' probability}

Another notion we'll use a great deal is \textbf{conditional} probability,
written with a vertical bar like P(red $| \overline{\textrm{female}}$). This
means ``if I know for a fact that the passenger is not female, what's the
probability that they're red-headed?'' Notice this is very different than the
joint probability. With conditional probability, I'm assuming that the second
event \textit{does} occur, and then asking how that revises my original
prediction. Unlike with joint probability, adding a vertical bar \textit{can}
make the answer increase.

Consider this example: what's the probability that a random U.S. citizen has a
Pinterest account? Let's estimate that as P(pinterest) = .15. But what if we
learned that the chosen person was female? What is P(pinterest $|$ female)? Just
to shamelessly stereotype, I'm going to go out on a limb and say it's higher;
perhaps .26. On the other hand, what if we learned the person was male? What
is P(pinterest $|$ male)? I'll say it's more like .04. In both cases, we
adjusted our estimated probability in light of additional information: up in
one case, and down in the other.

This is the heart of what's called ``\textbf{Bayesian reasoning}'' (you'll see
why in the next section) which is where our classifier gets its name.

Two other terms that come up frequently are ``\textbf{prior probability}'' and
``\textbf{posterior probability}''. The ``\textbf{prior}'' is the probability of
an event \textit{without} applying conditioning. In the previous example .15
was the ``prior probability'' that the person was a pinterest user. When we
learned our subject was female, we revised that to the ``\textbf{posterior}''
of .26.

Finally, note that P(A$|$B) is \textbf{\textit{not}} the same thing as
P(B$|$A). In general, these are very different numbers. Work out the pinterest
example further in your head to see why.

\subsection{The relationship between joint and conditional}

Mathematically, here's how these two concepts are related. If A and B are two
events, then:

\vspace{-.2in}
\begin{align*}
P(A|B) = \frac{P(A,B)}{P(B)}
\end{align*}

Suppose I knew that the probability of a random U.S. citizen being a
\textit{male pinterest user} was .04. And now suppose I knew that about half
of U.S. citizens are male. Then, I could compute:

\vspace{-.2in}
\begin{align*}
P(\textrm{male,pinterest}) = .04 \\
P(\textrm{male}) = .5 \\
P(\textrm{pinterest} | \textrm{male}) &= \frac{P(\textrm{male},
\textrm{pinterest})}{P(\textrm{pinterest})} \\
                    &= \frac{P(\textrm{pinterest},
\textrm{male})}{P(\textrm{pinterest})} \\
                    &= \frac{.04}{.15} = .267 \\
\end{align*}

\subsection{Bayes' Rule}

Knowing how joint and conditional probabilities are related, and knowing that
P(A,B) = P(B,A), we can derive the following law:

\vspace{-.2in}
\begin{align*}
P(A|B) = \frac{P(B|A) \cdot P(A)}{P(B)}
\end{align*}

This is called \textbf{Bayes' Rule} (sometimes \textbf{Bayes' Law} or
\textbf{Bayes' Theorem}) after the 18th-century Presbyterian minister Thomas
Bayes. It is one of the most influential (and controversial, actually)
statements in statistics and even all of mathematics.

\subsection{Independence}

We need one more background concept before we attack Na\"{i}ve Bayes Classifiers,
and that is the \textbf{independence} of events.

Simply put, two events are ``independent'' if \textit{whether one of them
occurs has no bearing on whether the other one occurs}. It does \textit{not}
mean the events are mutually exclusive. It means that the probability of each
one of them occurring has nothing to do with the other.

It's best seen through an example. Suppose I chose a random U.S. citizen, and
I observe two things: whether or not they're white, and whether or not they're
female. If I asked you for your estimate of the ``prior'' for the race
question, you'd probably come up with something like P(white) = .6. Now I ask
you: suppose, before observing their race, that you learned your subject was
female. Does this cause you to update your ``prior'' estimate to a more
accurate ``posterior''?

The answer is no. Whether or not someone is white is \textit{independent} of
whether they're female. You thought P(white) was .6, and you still think
P(white$|$female) is .6. The opposite was true in the pinterest example: your
estimate of their pinterest status \textit{did} change when you learned about
their gender. Therefore, those two events are \textit{not} independent.

Both of the following mathematical identities are true \textit{\textbf{only
\underline{if}}} the events $A$ and $B$ are independent:
\vspace{-.2in}
\begin{align*}
P(A,B) &= P(A) \cdot P(B) \\
P(A|B) &= P(A)
\end{align*}

\section{The Na\"{i}ve Bayes algorithm}

Suppose we're designing a security system for a Galactic Cruiser. One thing we
need to do is identify possible evildoers on our security cameras. Given what
we can observe about them, can we tell good guys from bad guys?

Here's some training data:
\begin{Verbatim}[fontsize=\small,samepage=true,frame=none]
    beard  height  demeanor lightsaber  type
       no   short  menacing      green  Sith
      yes  medium      calm      green  Jedi
       no    tall  menacing        red  Sith
      yes  medium      calm       blue  Jedi
      yes  medium      calm       blue  Jedi
       no  medium      calm       blue  Sith
       no    tall  menacing        red  Jedi
       no  medium  menacing        red  Sith
       no   short      calm       blue  Jedi
      yes  medium  menacing      green  Sith
                    . . .
\end{Verbatim}

Our features are the first four columns: whether the person in question has a
\texttt{beard}, what their \texttt{height} is (clustered in three broad
ranges), whether their \texttt{demeanor} is \texttt{menacing} or
\texttt{calm}, and what color \texttt{lightsaber} they wield. The final
column is our label: this person is known to be a \texttt{Jedi} (good guy) or
a \texttt{Sith} (bad guy).

A test point that we might want to evaluate looks like this:
\begin{Verbatim}[fontsize=\small,samepage=true,frame=none]
    beard  height  demeanor lightsaber  type
      yes   short  menacing        red   ?? 
\end{Verbatim}

A short, menacing, bearded person with a red lightsaber has just showed up on
security camera \#305. Should we sound the alarm? That, of course, depends
upon their (as yet unknown) \texttt{type}. We want to predict this based on
our past examples.

The \textbf{Na\"{i}ve Bayes} classification algorithm applies conditional
probability to this problem. Really, what we're asking is: ``which of
the following two quantities is greater:

\begin{enumerate}
\itemsep.05em
\item P(Jedi $|$ yes, short, menacing, red),  or
\item P(Sith $|$ yes, short, menacing, red)?''
\end{enumerate}
 
Note that we're combining joint and conditional probabilities here. Given that
we have the \textit{joint} event ``bearded, short, menacing, and red,'' what's
the \textit{conditional} probability of them being a Jedi (or Sith)? If we
estimate the first as .8 and the second as .2, we won't sound the alarm.

\section{Using Bayes' rule}

With Bayes' rule, we can rewrite these as:

\vspace{-.2in}
\begin{align}
P(\textrm{Jedi} \ | \ \textrm{yes}, \textrm{short}, \textrm{menacing},
\textrm{red}) &=
\frac{P(\textrm{yes}, \textrm{short}, \textrm{menacing}, \textrm{red} \ | \ 
\textrm{Jedi}) \cdot P(\textrm{Jedi})}{P(\textrm{yes}, \textrm{short},
\textrm{menacing},
\textrm{red})} \label{eqn:jedi} \\
P(\textrm{Sith} \ | \ \textrm{yes}, \textrm{short}, \textrm{menacing},
\textrm{red}) &=
\frac{P(\textrm{yes}, \textrm{short}, \textrm{menacing}, \textrm{red} \ | \ 
\textrm{Sith}) \cdot P(\textrm{Sith})}{P(\textrm{yes}, \textrm{short},
\textrm{menacing},
\textrm{red})}  \label{eqn:sith}
\end{align}

Let's focus on the first of these two expressions, since the second is
identical except for the word Sith.

We've turned things on their head. In essence, instead of asking ``given these
physical features, are they a Jedi?'' we're asking ``given they're a Jedi,
would they have these physical features?'' 

The right-hand side of that equation has three parts. Ignore the denominator
for a moment. The right-part of the numerator, P(Jedi), is our \textbf{prior}
probability. That's easy to estimate from the training data: we just count up
what percentage of our samples are Jedi overall.

The left-hand side of the numerator needs more work. It reads:

\vspace{-.2in}
\begin{align*}
P(\textrm{yes}, \textrm{short}, \textrm{menacing}, \textrm{red} \ | \ \textrm{Jedi}) 
\end{align*}

Now \textbf{\textit{\underline{if}}} it just so happened that those four
features were \textit{independent} of one another, this would actually be
super easy to compute. We could just do this:

\vspace{-.2in}
\begin{align*}
P(\textrm{yes} \ |\ \textrm{Jedi}) \cdot
P(\textrm{short} \ |\ \textrm{Jedi}) \cdot
P(\textrm{menacing} \ |\ \textrm{Jedi}) \cdot
P(\textrm{red} \ |\ \textrm{Jedi})
\end{align*}

Under this assumption, we've split the expression into \textit{each feature
individually.} These are all now trivially obtainable from the training data.
What's P(yes $|$ Jedi)? Look at \textit{only the Jedi} in the training data,
and count up how many ``yesses'' and ``nos'' there are. The fraction that are
yesses is our estimate of P(yes $|$ Jedi).\footnote{If you have trouble seeing
this, go back to the definition of conditional probability. We're really
calculating $\dfrac{P(\textrm{yes},\textrm{Jedi})}{P(\textrm{Jedi})}$ here,
which is ``the number of bearded Jedi divided by the number of total Jedi.''}

This is called \textbf{the Na\"{i}ve Bayes assumption}. Of course it's probably
not true that beardedness, height, demeanor, and lightsaber color are all
independent of each other; for one thing, if some or all of these features are
indicative of Jedi-ness (which we hope), then they certainly \textit{do} vary
together. But the Na\"{i}ve Bayes classifier says ``aw heck, let's just throw
caution to the wind and see how things work out.''

\section{A numerical example}

Let's work out the numerical answer (numerator only) for the training data
listed earlier. First, we consider \textit{only} the five \texttt{Jedi} rows,
and obtain for equation (\ref{eqn:jedi}):

\vspace{-.2in}
\begin{align*}
P(\textrm{yes} \ |\ \textrm{Jedi}) \cdot
P(\textrm{short} \ |\ \textrm{Jedi}) \cdot
P(\textrm{menacing} \ |\ \textrm{Jedi}) \cdot
P(\textrm{red} \ |\ \textrm{Jedi}) \cdot P(\textrm{Jedi}) = \\
\frac{3}{5} \cdot \frac{1}{5} \cdot \frac{1}{5} \cdot \frac{1}{5} \cdot
\frac{5}{10} = .0024\\
\end{align*}

Now let's do the same thing for equation (\ref{eqn:sith}) and the five \texttt{Sith} rows:

\vspace{-.2in}
\begin{align*}
P(\textrm{yes} \ |\ \textrm{Sith}) \cdot
P(\textrm{short} \ |\ \textrm{Sith}) \cdot
P(\textrm{menacing} \ |\ \textrm{Sith}) \cdot
P(\textrm{red} \ |\ \textrm{Sith}) \cdot P(\textrm{Sith}) = \\
\frac{1}{5} \cdot \frac{1}{5} \cdot \frac{4}{5} \cdot \frac{2}{5} \cdot
\frac{5}{10} = .0064\\
\end{align*}

\subsection{Throw away the denominator}

These are the two numerators. To get true probabilities, we have to divide by
the denominators, of course. But the nice thing is: \textit{the denominators
are the \underline{same} in the two cases!} For both expressions, we would
need to divide by P(\textrm{yes}, \textrm{short}, \textrm{menacing},
\textrm{red}) (in words, ``what's the probability that any random person would
be a bearded, short, menacing person with a red lightsaber?'') which of course
has the same value even if we compute it twice. For this reason, we
\textit{throw away the denominator} since in the end all we really care about
is ``which is higher, P(Jedi $|$ stuff) or P(Sith $|$ stuff)?'' In this case,
it's P(Sith $|$ stuff) and so we boldly predict: ``this person looks like a
Sith: sound the alarm.''

\section{A Python implementation}

How can we automate this in Python? Let's assume that our data is in a Pandas
DataFrame, as shown above. Now imagine writing this function:

\begin{Verbatim}[fontsize=\small,samepage=true,frame=single]
def predict(beard, height, demeanor, lightsaber):
    ...
\end{Verbatim}

This function takes as input a test (unlabeled) point and returns a prediction
for it: either the word ``Jedi'' or the word ``Sith.''


We'll do most of the work outside the function, actually. We'll read in the
file and then pre-compute all the probabilities that we'll need for our
calculations. If you think it through, the calculations above imply the need
for the following probabilities:

\pagebreak
\begin{itemize}
\itemsep.1em
\item $P(label)$ for every one of the target labels (these are the ``priors'')
\item $P(value|label)$ for every value in a column and every target label (the
conditionals)
\end{itemize}

In our case, this amounts to:

\vspace{.1in}
\begin{compactitem}
\item Priors:
    \begin{compactitem}
    \item $P(\textrm{Jedi})$
    \item $P(\textrm{Sith})$
    \end{compactitem}

\item Conditionals:
    \begin{compactitem}
    \item $P(\textrm{no}|\textrm{Jedi})$ and
$P(\textrm{no}|\textrm{Sith})$\footnote{Recall that ``\texttt{yes}'' and
``\texttt{no}'' refer to whether the person in question has a beard.}
    \item $P(\textrm{yes}|\textrm{Jedi})$ and $P(\textrm{yes}|\textrm{Sith})$
    \item $P(\textrm{tall}|\textrm{Jedi})$ and $P(\textrm{tall}|\textrm{Sith})$
    \item $P(\textrm{medium}|\textrm{Jedi})$ and $P(\textrm{medium}|\textrm{Sith})$
    \item $P(\textrm{short}|\textrm{Jedi})$ and $P(\textrm{short}|\textrm{Sith})$
    \item $P(\textrm{menacing}|\textrm{Jedi})$ and $P(\textrm{menacing}|\textrm{Sith})$
    \item $P(\textrm{calm}|\textrm{Jedi})$ and $P(\textrm{calm}|\textrm{Sith})$
    \item $P(\textrm{red}|\textrm{Jedi})$ and $P(\textrm{red}|\textrm{Sith})$
    \item $P(\textrm{green}|\textrm{Jedi})$ and $P(\textrm{green}|\textrm{Sith})$
    \item $P(\textrm{blue}|\textrm{Jedi})$ and $P(\textrm{blue}|\textrm{Sith})$
    \end{compactitem}

\end{compactitem}

These are all straightforward to compute with Pandas. First, we read the file
into a DataFrame (called ``\texttt{w}'' for ``warriors''):

\begin{Verbatim}[fontsize=\small,samepage=true,frame=single]
import numpy as np
import pandas as pd
w = pd.read_csv("warriors.csv")
\end{Verbatim}

I find it convenient to create a dictionary, called something like
``\texttt{probs}'' to store all the probabilities. To compute the priors, all
we have to do is figure out what fraction of the rows go with each label:

\begin{Verbatim}[fontsize=\small,samepage=true,frame=single]
probs = {}
probs['Jedi'] = len(w[w.type=='Jedi'])/len(w)
probs['Sith'] = len(w[w.type=='Sith'])/len(w)
\end{Verbatim}

For the conditionals like ``$P(\textrm{calm}|\textrm{Sith})$'', we ask:
\textit{considering only the rows that have \texttt{Sith} labels}, what
fraction also have \texttt{calm} labels? That boils down to this:

\footnotesize
\begin{Verbatim}[fontsize=\small,samepage=true,frame=single]
probs['calm|Sith'] = (len(w[(w.type=='Sith') & (w.demeanor=='calm')]) /
    len(w[(w.type=='Sith')]))
\end{Verbatim}
\normalsize

The other 9 conditionals follow exactly the same pattern. (It may strike you
that we're repeating a lot of very similar code here, and making slight edits
to each one, and that there's a better, more general way. You're right. Stay
tuned for the next section.)


Now for the actual \texttt{predict()} function. All we have to do is multiply
everything out per the above equations and see which of the two numerators is
greater. The result is in Figure~\ref{fig:predictFunc}. Easy peasy. Note that
we don't use the \texttt{w} variable inside this function at all. We've
already computed all the \texttt{probs}, and once we've done that, the
DataFrame and the data itself ceases to matter. We run the function just by
passing it the values derived from our spycam:

\begin{figure}
\begin{Verbatim}[fontsize=\small,samepage=true,frame=single]
def predict(beard, height, demeanor, lightsaber):
    jediness = (probs[beard+"|Jedi"] *
                probs[height+"|Jedi"] *
                probs[demeanor+"|Jedi"] *
                probs[lightsaber+"|Jedi"] *
                probs["Jedi"])
    sithness = (probs[beard+"|Sith"] *
                probs[height+"|Sith"] *
                probs[demeanor+"|Sith"] *
                probs[lightsaber+"|Sith"] *
                probs["Sith"])
    if jediness > sithness:
        return "Jedi"
    else:
        return "Sith"
\end{Verbatim}
\caption{Our Na\"{i}ve Bayes ML \texttt{predict} function.}
\label{fig:predictFunc}
\end{figure}

\begin{Verbatim}[fontsize=\small,samepage=true,frame=single]
>>> predict("yes","tall","calm","blue")
    'Jedi'
>>> predict("no","tall","menacing","red")
    'Sith'
\end{Verbatim}


\section{Estimating confidence}

Our \texttt{predict()} function, as written, just gives a binary answer:
either \texttt{Sith} or \texttt{Jedi}. It may be useful to know \textit{how
confident} it is in its prediction for any given input. All we need to do is
compute the correct ratio. Remember that \texttt{jediness} and
\texttt{sithness} are just the \textit{numerators} of probabilities, and that
those probabilities have the same denominator. Therefore, in order to compute
the true $P(\textrm{Jedi}|e)$, where $e$ is all our evidence (calm, no beard,
and so forth), we just need to compute:

\begin{align*}
P(\textrm{Jedi}|e) = \frac{\textrm{jediness}}{\textrm{jediness} + \textrm{sithness}}
\end{align*}

and similarly, 

\begin{align*}
P(\textrm{Sith}|e) = \frac{\textrm{sithness}}{\textrm{jediness} + \textrm{sithness}}
\end{align*}

The denominator drops out when we divide the quantities anyway.

So let's enhance our \texttt{predict()} function to return a list of
\textit{two} values: the prediction, and the confidence level (on a scale of
.5 to 1, meaning ``only 50\% confident in our prediction'' all the way to
``absolutely 100\% confident in our prediction.'' To wit:

\begin{Verbatim}[fontsize=\small,samepage=true,frame=single]
def predict(beard, height, demeanor, lightsaber):
    ...
    if jediness > sithness:
        return ["Jedi", jediness/(jediness+sithness)]
    else:
        return ["Sith", sithness/(jediness+sithness)]
\end{Verbatim}

Taking it for a spin:

\begin{Verbatim}[fontsize=\small,samepage=true,frame=single]
>>> predict("no","tall","menacing","red")
    ['Sith', 0.9760957241941071]
>>> predict("yes","tall","menacing","blue")
    ['Jedi', 0.7298292697172792]
\end{Verbatim}

We're incredibly confident that our beardless, tall, menacing, red-wielder is
a Sith, but not nearly so confident that our bearded, tall, menacing
blue-wielder is a Jedi.

\vspace{.2in}
\section{Coding it more generally}

The above calculation of the probabilities left something to be desired, since
it was hardcoded and repetitive. We can improve on this by writing our code in
a more general way, rather than explicitly referring to problem-specific names
like \texttt{'Jedi'} and \texttt{'demeanor'} and \texttt{'blue'}. Here's one
way to do all the calculations:


\begin{Verbatim}[fontsize=\small,samepage=true,frame=single]
df = pd.read_csv("...")

probs = {}
features = df.columns[:-1]
target = df.columns[-1]

for label in df[target].drop_duplicates():
    probs[label] = len(df[df[target] == label])/len(df)
    for feature in features:
        for value in df[feature].drop_duplicates():
            probs[value+"|"+label] = (
                len(df[(df[target] == label) & (df[feature] == value)]) / 
                    len(df[df[target] == label]))
\end{Verbatim}

Here's we're assuming (1) the last column of the DataFrame is the target, and
(2) all the values in the columns are unique (\textit{e.g.}, that we won't
have an ``\texttt{out-of-state}'' column with some ``\texttt{yes}'' values and
also have a ``\texttt{married}'' column with ``\texttt{yes}'' values. But those
are pretty much our only assumptions.

The \texttt{predict()} method itself will be generalized as follows: instead
of taking exactly four parameters named ``\texttt{beard}'',
``\texttt{demeanor}'', \textit{etc.}, it will take a \textit{list} of
parameter values, so that it can be called like this:

\begin{Verbatim}[fontsize=\small,samepage=true,frame=single]
>>> predict(['yes','tall','menacing','blue'])
\end{Verbatim}

Here's one implementation of the function:

\begin{Verbatim}[fontsize=\small,samepage=true,frame=single]
def predict(values):
    scores = []
    top_prob = 0
    top_answer = ""
    for label in df[target].drop_duplicates():
        prob = probs[label]                     # Start with the prior
        for value in values:
            prob *= probs[value+"|"+label]
        scores.append(prob)
        if prob > top_prob:
            top_prob = prob
            top_answer = label
    return [ top_answer, top_prob / sum(scores) ]
\end{Verbatim}

We use the variable \texttt{scores} to keep a list of the numerators for each
target label: this is only necessary to compute the confidence level in the
result.

\section{Measuring performance}

So how well does this classifier work? This depends on how we measure its
``goodness.'' Commonly, we'll simply ask, ``what percentage of the time is it
right?'' and count errors the same in either direction (calling ``Sith'' for an
actual Jedi, or calling ``Jedi'' for an actual Sith).

As we know, we can't test a classifier's performance on the data it's trained
on and expect to get anything accurate. So we'll need to split our data into a
training set and a test set. This can be done easily in Pandas with:

\begin{Verbatim}[fontsize=\small,samepage=true,frame=single]
complete_data_set = pd.read_csv("warriors.csv")
train = complete_data_set.sample(frac=.7)
test = complete_data_set.drop(train.index)
w = train
\end{Verbatim}

This call to \texttt{.sample()} chooses 70\% of the original, complete
DataFrame's rows and copies them to a DataFrame called ``\texttt{train}''. Then
the next line makes ``\texttt{test}'' a DataFrame of all the others (the
remaining 30\%). For convenience, we then copy \texttt{train} into the
\texttt{w} variable so that none of our previous code has to change.

All our probabilities are thus computed based \textit{only} on the training
data (the random 70\% of the data set). We can then use the \texttt{test} set
to compute our accuracy. All we need to do is iterate through the
\texttt{test} set, passing each row to \texttt{predict()} and seeing whether
we get the correct answer:

\begin{Verbatim}[fontsize=\footnotesize,samepage=true,frame=single]
num_correct = 0
for i in range(len(test)):
    row = test.iloc[i]
    p = predict(row['beard'],row['height'],row['demeanor'],row['lightsaber'])
    if p[0] == row['type']:
        num_correct += 1

print("We got " + str(num_correct/len(test)*100) +  
    "% correct on the test data.")
\end{Verbatim}

We get a different answer every time we run our program, since it will choose
a different training/test split in each case.

\begin{Verbatim}[fontsize=\small,samepage=true,frame=single]
We got 82.63% correct on the test data.
\end{Verbatim}

\section{Numerical/continuous variables}

In all the examples above, we had categorical features. Computing something
like ``$P(\textrm{calm}|\textrm{Jedi})$'' is thus just a matter of counting the
\texttt{Jedi} rows which (exactly) match the value \texttt{calm}. But what if
we have a numeric feature, like a person's age or salary?

One quick-and-dirty way to deal with numeric features is to ``bin'' them, kind
of like a histogram. Note that this is actually what we already did with the
``\texttt{height}'' column in the Jedi example: instead of heights like 4'9"
and 6'6", we had values like \texttt{short} and \texttt{tall}.

This has obvious drawbacks. For one, we're throwing away all the detail
present in the measurement and treating different values exactly alike. For
another, it's up to us to somewhat arbitrarily choose how many bins to use and
where the demarcation lines are. This has an \textit{ad hoc} feel to it.

Better is to treat the continuous variable as what it really is: a number.
\textbf{Gaussian Na\"{i}ve Bayes} is probably the most common way to deal with
this in a principled fashion: we make the assumption that the numeric feature
is distributed \textbf{normally} (``normal''==``Gaussian''), and estimate its
mean and standard deviation from the training data. Then, to estimate
$P(value|target\_label)$ for any particular test data point's value, we choose
the estimated value of its \textbf{probability density} and use that.

Suppose we had a numeric \texttt{height} feature (measured in inches) instead
of a categorical one. Our data might look like this:

\begin{Verbatim}[fontsize=\small,samepage=true,frame=none]
    beard  height  demeanor lightsaber  type
       no    62.5  menacing      green  Sith
      yes    68.3      calm      green  Jedi
       no    78.1  menacing        red  Sith
      yes    63.9      calm       blue  Jedi
      yes    61.9      calm       blue  Jedi
       no    68.2      calm       blue  Sith
       no    80.1  menacing        red  Jedi
       no    33.5      calm      green  Jedi
                    . . .
\end{Verbatim}

Our Gaussian Na\"{i}ve Bayes assumption will be that \textit{for each class, the
\texttt{height} feature is distributed normally}. If we take all the
\texttt{Jedi} rows, and calculate their \texttt{height}s' mean and standard
deviation, we get:

\begin{Verbatim}[fontsize=\small,samepage=true,frame=single]
>>> warriors[warriors.type=='Jedi'].height.mean()
    64.0571428571429
>>> warriors[warriors.type=='Jedi'].height.std()
    9.661931704216865
\end{Verbatim}

Performing the same calculation for the \texttt{Sith} rows gives us:

\begin{Verbatim}[fontsize=\small,samepage=true,frame=single]
>>> warriors[warriors.type=='Sith'].height.mean()
    68.3902173913043
>>> warriors[warriors.type=='Sith'].height.std()
    9.319673312968604
\end{Verbatim}

Looks like Sith Lords are a bit taller than Jedi in general, and slightly less
variable.

\vspace{.1in}
Let's say we're doing a prediction for a data point with a height of 69
inches. We'll need to know the value of the density at $x$=69 for both our
estimated Jedi and estimated Sith distributions. Using the SciPy library, we
can get ``frozen'' density functions for each distribution:

\begin{Verbatim}[fontsize=\small,samepage=true,frame=single]
import scipy.stats
jedi_height = scipy.stats.norm(64.0571428571429, 9.661931704216865)
sith_height = scipy.stats.norm(68.3902173913043, 9.319673312968604)
\end{Verbatim}

and then, to calculate the densities, we just call \texttt{.pdf()} on those
objects:

\begin{Verbatim}[fontsize=\small,samepage=true,frame=single]
>>> jedi_height.pdf(69)
    0.036230663486274517
>>> sith_height.pdf(69)
    0.042717943988373371
\end{Verbatim}

Not surprisingly, the Sith number comes out a fair bit higher, since according
to our estimated density, the average Sith's height is taller than the average
Jedi's.

Keep in mind that the ``.0362'' and ``.0427'' numbers here are a bit hard to
interpret. They are \textit{not} ``the probability that a Jedi/Sith is exactly
69 inches tall'' -- that probability, as with all continuous random variables,
is \textit{zero}. Instead, it's ``the value of the density at $x$=69'' which
isn't really in any units that make sense. It's okay, though, since we are
throwing away out denominator and doing relative comparisons anyway.

By the way, there are also other useful things you can call on a frozen SciPy
distribution, including \texttt{.cdf()} to get the value of the cumulative
density function at any point, \texttt{.ppf()} to get quantile,
\texttt{.rvs()} to generate random variates, ordinary statistical things like
\texttt{.median()}, \texttt{.mean()}, \texttt{.std()}, and other more exotic
statistical things like moments, confidence intervals, differential entropy,
survival functions, \textit{etc.} See the documentation for details.

\subsection{Incorporating into the classifier}

Let's add this feature to our classifier. First, the not-so-generally-coded
version, which had lines like this:

\footnotesize
\begin{Verbatim}[fontsize=\small,samepage=true,frame=single]
probs['tall|Sith'] = (len(w[(w.type=='Sith') & (w.demeanor=='tall')]) /
    len(w[(w.type=='Sith')]))
\end{Verbatim}
\normalsize

Since \texttt{height} is now numeric, we need to replace all those
height-related lines with code like this:

\begin{Verbatim}[fontsize=\small,samepage=true,frame=single]
sith_heights = w[w.type == 'Sith'].height
sith_height = scipy.stats.norm(sith_heights.mean(),sith_heights.std())
jedi_heights = w[w.type == 'Jedi'].height
jedi_height = scipy.stats.norm(jedi_heights.mean(),jedi_heights.std())
\end{Verbatim}

``\texttt{sith\_height}'' is now our frozen distribution function. In 
\texttt{predict()}, then, we simply use it as another multiplier:

\begin{Verbatim}[fontsize=\small,samepage=true,frame=single]
def predict(beard, height, demeanor, lightsaber):
    jediness = (probs[beard+"|Jedi"] *
                probs[demeanor+"|Jedi"] *
                probs[lightsaber+"|Jedi"] *
                probs["Jedi"])
    sithness = (probs[beard+"|Sith"] *
                probs[demeanor+"|Sith"] *
                probs[lightsaber+"|Sith"] *
                probs["Sith"])

    jediness *= jedi_height.pdf(height)
    sithness *= sith_height.pdf(height)

    if jediness > sithness:
        return "Jedi"
    else:
        return "Sith"
\end{Verbatim}

Notice we've removed the ``\texttt{probs[height+"|Jedi"]}'' lines,
since they're no longer applicable: \texttt{height} is numeric now. Instead,
we multiply our \texttt{jediness}/\texttt{sithness} by the value of the
corresponding class's estimated density at that point.

To incorporate this numeric variable into the more general version of the
code, we need to intelligently distinguish between numeric and categorical
columns. This can be done via the \texttt{.dtype} accessor on a DataFrame
column:

\begin{Verbatim}[fontsize=\small,samepage=true,frame=single]
>>> print(w.lightsaber.dtype)
    dtype('O')
>>> print(w.height.dtype)
    dtype('float64')
\end{Verbatim}

(The ``\texttt{O}'' stands for ``object,'' which is Pandas' non-intuitive way of
telling you that the column is something non-numeric, like a string.) You can
directly compare the result to a value like ``\texttt{np.dtype("float64")}'',
but the safer and preferred way is to use NumPy's \texttt{.issubdtype()}
function which more intelligently detects any numeric type (ints, floats,
\textit{etc.}, no matter the precision).

Our generic implementation begins with:

\begin{Verbatim}[fontsize=\small,samepage=true,frame=single]
probs = {}
densities = {}
target = df.columns[-1]
numeric_features = [ c for c in df.columns[:-1] 
    if np.issubdtype(df[c].dtype, np.number) ]
categorical_features = [ c for c in df.columns[:-1]
    if not np.issubdtype(df[c].dtype, np.number) ]
\end{Verbatim}

We're creating a \texttt{densities} dictionary, somewhat similar to our
\texttt{probs} dictionary, to store feature-specific information for our
numeric variables (\texttt{probs} will continue to store information for
categorical variables.) Instead of simple numeric values, \texttt{densities}
will store frozen density functions.

We now populate \texttt{densities} (and \texttt{probs}):

\begin{Verbatim}[fontsize=\small,samepage=true,frame=single]
for label in df[target].drop_duplicates():
    rows_with_this_label = df[df[target] == label]
    probs[label] = len(rows_with_this_label)/len(df)
    for feature in numeric_features:
        density = scipy.stats.norm(rows_with_this_label[feature].mean(),
            rows_with_this_label[feature].std())
        densities[feature+"|"+label] = density
    for feature in categorical_features:
        for value in df[feature].drop_duplicates():
            probs[value+"|"+label] = (
                len(df[(df[target] == label) & (df[feature] == value)]) / 
                    len(df[df[target] == label]))
\end{Verbatim}

Our generic \texttt{predict()} function needs to be adjusted somewhat, since
if we just pass a list of values, we can't tell \textit{which} feature a
numeric value is supposed to go with. There are various ways to solve this,
but we'll choose to pass a key/value pair string here (like
\texttt{"height=69.3"}):

\begin{Verbatim}[fontsize=\footnotesize,samepage=true,frame=single]
def predict(values):
    scores = []
    top_prob = 0
    top_answer = ""
    for label in df[target].drop_duplicates():
        prob = probs[label]    # Start with the prior
        for value in values:
            if "=" in value:
                # Numeric feature.
                numeric_feature, numeric_value = value.split("=")
                numeric_value = float(numeric_value)
                prob *= densities[numeric_feature+"|"+label].pdf(numeric_value)
            else:
                # Categorical feature.
                prob *= probs[value+"|"+label]
        scores.append(prob)
        if prob > top_prob:
            top_prob = prob
            top_answer = label
    return [ top_answer, top_prob / sum(scores) ]
\end{Verbatim}

And we call it like this:
\begin{Verbatim}[fontsize=\small,samepage=true,frame=single]
>>> predict(['yes','height=65','menacing','blue'])
    ['Jedi', 0.85863336769968801]
>>> predict(['yes','height=85','menacing','blue'])
    ['Jedi', 0.72742228657164021]
\end{Verbatim}

How's this thing do?

\begin{Verbatim}[fontsize=\small,samepage=true,frame=single]
We got 82.33% correct on the test data.
\end{Verbatim}

Not too shabby.

